% ---
% id: foundations/def:domination
% kind: def
% title: "Domination of one operator family by another at block weights"
% summary: >-
%   A family G dominates a family F at block weights q when G − (1−q)F is
%   positive semidefinite for every outcome on every block, one weight per
%   block. The residual family is what G holds beyond the weighted F,
%   rescaled by the weight.
% prominence: standard
% status: accepted
% uses: [def:sector-map]
% mentions: [meat-core/prop:optol-nobound]
% assumes: []
% refs:
%   - "nahar2025imperfect, Thm. 3, Eq. (20)"
% provenance: newly-derived
% notes: "The relation is the hypothesis of the generic noise channel of nahar2025imperfect, Thm. 3, Eq. (20), with G in the role of the noisy POVM and F in that of the noiseless one; the channel itself is not stated in this module. No requirement uses the relation: a snap is not related to its model by domination, and a declared device weight is future scope (spec/future-scope.md section 13)."
% lean:
% ---

Two \tref{exactly admissible} measurements that are entrywise close can have key-rate
bounds far apart, by an amount the distance between them does not control
(\xref{meat-core/prop:optol-nobound}). A small change can alter which
outcomes a measurement treats as impossible on some direction, and a bound on
its size says nothing about that. Domination is a relation between two families
under which the dominated family gains no support; it bounds the dominated
family from above and not from below. It is the hypothesis under which a measurement $G$ is
a noise channel followed by $F$ together with classical flags of weight at
most $q_s$ on block $s$ \cite[Thm.~3, Eq.~(20)]{nahar2025imperfect}.

\begin{definition}\label{def:domination}
Let $\mathfrak{m}$ be a \tref{sector map} with blocks $s\in\mathsf{S}$ of
dimensions $d_s$ (Definition~\ref{def:sector-map}), let $J$ be a finite non-empty
alphabet, and let $G=\{G_j\}_{j\in J}$ and $F=\{F_j\}_{j\in J}$ be families of
\tref[block list]{block lists} against $\mathfrak{m}$ whose blocks are
Hermitian and positive semidefinite; their entries are arbitrary complex
numbers. A family is \emph{complete at} a real number $c$ on block $s$ when
$\sum_{j\in J}G^{(s)}_j=c\,\idon{d_s}$. For weights $q=(q_s)_{s\in\mathsf{S}}$
with $q_s\in[0,1)$, the family $G$
\term[said of an operator family G and an operator family F over the same
alphabet and sector map, at block weights q in [0, 1): G − (1−q)F is positive
semidefinite for every outcome on every block, with the weight of that
block]{dominates} $F$ at $q$ when
\begin{equation}\label{eq:domination}
  G^{(s)}_j-(1-q_s)\,F^{(s)}_j\;\succeq\;0
  \qquad\text{for every }j\in J\text{ and }s\in\mathsf{S}.
\end{equation}
The weight $q_s$ belongs to the block $s$ and is the same for every outcome.
On a block $s$, the
\term[on a block with weight q, the family (G − (1−q)F)/q where q is positive,
and G itself where q is zero]{residual family} of the domination is
$R^{(s)}_j\coloneqq q_s^{-1}\big(G^{(s)}_j-(1-q_s)F^{(s)}_j\big)$ for
$j\in J$ where $q_s>0$, and $R^{(s)}_j\coloneqq G^{(s)}_j$ where $q_s=0$.
For a Hermitian matrix $X$ we write $\mathrm{supp}\,X$ for its range, the
orthogonal complement of its kernel.
\end{definition}
